\documentclass[a4paper,12pt]{article}
\usepackage[margin=2.5cm]{geometry}
\usepackage{graphicx}
\usepackage{booktabs}
\usepackage{amsmath}
\usepackage{amssymb}
\usepackage{caption}
\usepackage{subcaption}
\usepackage{hyperref}
\usepackage{xcolor}
\usepackage{tabularx}

\title{\textbf{Result 3: Network Throughput vs Packet Generation Rate}\\
\large ACLIF --- Autonomous Cross-Layer Intelligence Framework}
\author{B.~Tharuni~Sri~Sai\\
\small School of Computer Science and Engineering, VIT-AP University, Amaravati, India\\
\small ORCID: 0009-0004-9561-8985}
\date{\today}

\begin{document}
\maketitle
\tableofcontents
\newpage

\section{Experiment Overview}
This document presents the detailed analysis for \textbf{Result 3}: Network Throughput vs Packet Generation Rate,
as part of the ACLIF cross-layer IoT optimization framework evaluation.

\section{Result Figure}
\begin{figure}[htbp]
  \centering
  \includegraphics[width=0.92\textwidth]{../figures/fig_throughput_rate.png}
  \caption{Network Throughput vs Packet Generation Rate. Results obtained from NS3 v3.38 simulations over a
  100-node heterogeneous IoT network ($200\times200$~m$^2$), averaged over 10 runs
  with 95\% confidence interval shading.}
  \label{fig:throughput_vs_rate}
\end{figure}

\section{Analysis and Discussion}
Network throughput is measured against offered load from 1 to 10 pkt/s. ACLIF's
transport rate adaptation prevents congestion collapse at high loads by proactively
reducing send rates before queues overflow. At $\lambda=10$ pkt/s, ACLIF achieves
83.7~kbps vs 69.8~kbps for IPSO (+19.9\%) and 51.6~kbps for FACSO (+62.2\%).
The DQN learns to predict queue overflow from RTT measurements in the transport
state sub-vector and preemptively engages rate backoff, a capability absent from
all reactive baselines.

\subsubsection*{Throughput Definition}
\begin{equation}
\Theta(t) = \frac{\sum_{i=1}^{N} \mathbf{1}[\mathrm{pkt}_i \text{ delivered at } t]}{\Delta t}
\end{equation}

\subsubsection*{Summary Statistics}
\begin{tabular}{lcc}
\toprule
Method & Throughput @ $\lambda{=}10$ (kbps) & \% vs ACLIF \\
\midrule
E-GLBR & 58.2  & $-$30.5\% \\
FDRL   & 67.4  & $-$19.5\% \\
FACSO  & 51.6  & $-$38.3\% \\
EEHCT  & 63.1  & $-$24.6\% \\
IPSO   & 69.8  & $-$16.6\% \\
\textbf{ACLIF} & \textbf{83.7} & --- \\
\bottomrule
\end{tabular}

\section{Experimental Configuration}
\begin{tabular}{ll}
\toprule
\textbf{Parameter} & \textbf{Value} \\
\midrule
Simulator      & NS3 v3.38 \\
Nodes          & 100 sensor nodes \\
Area           & $200 \times 200$~m$^2$ \\
PHY standard   & IEEE 802.15.4 \\
MAC protocol   & CSMA/CA \\
Packet payload & 512 bytes \\
Packet rate    & 5 pkt/s (default) \\
Runs per config & 10 (different seeds) \\
CI             & 95\% Student-$t$ \\
\bottomrule
\end{tabular}

\section{Baselines}
ACLIF is compared against: E-GLBR \cite{Benelhouri2023}, FDRL \cite{Suresh2024},
FACSO \cite{Yang2024204}, EEHCT \cite{Chaurasiya202325941}, and IPSO \cite{Zhang2025}.

\begin{thebibliography}{9}
\bibitem{Benelhouri2023} Benelhouri et al., ``E-GLBR Genetic Algorithm Routing,'' 2023.
\bibitem{Suresh2024} Suresh et al., ``FDRL Federated Deep RL Routing,'' 2024.
\bibitem{Yang2024204} Yang et al., ``FACSO Fuzzy Adaptive Cross-Layer Scheduling,'' 2024.
\bibitem{Chaurasiya202325941} Chaurasiya et al., ``EEHCT Energy-Efficient Clustering,'' 2023.
\bibitem{Zhang2025} Zhang et al., ``IPSO-Based Multi-Hop Routing,'' 2025.
\end{thebibliography}

\end{document}
